
% ============================================================
% FRENCH TENSES — FUTUR SIMPLE
% Version graphique LaTeX/TikZ
% Compilation : XeLaTeX
% ============================================================
\documentclass[a4paper,10pt]{article}

\usepackage[a4paper,top=9mm,bottom=8mm,left=11mm,right=11mm]{geometry}
\usepackage{fontspec}
\usepackage[french]{babel}
\usepackage{microtype}
\usepackage{tikz}
\usetikzlibrary{arrows.meta,calc,positioning}
\usepackage[most]{tcolorbox}
\usepackage{tabularx}
\usepackage{array}
\usepackage{xcolor}
\usepackage[normalem]{ulem}

\setmainfont{Noto Sans}
\setsansfont{Noto Sans}

% ------------------------------------------------------------
% PALETTE — FUTUR = ORANGE
% ------------------------------------------------------------
\definecolor{Orange}{HTML}{D8842F}
\definecolor{OrangeDark}{HTML}{9A5716}
\definecolor{OrangeLight}{HTML}{FFF1DF}
\definecolor{OrangeLine}{HTML}{E7B77E}
\definecolor{Border}{HTML}{E7C7A6}

\definecolor{Red}{HTML}{D52F35}
\definecolor{RedLight}{HTML}{FCEBED}
\definecolor{Gold}{HTML}{B27A00}
\definecolor{GoldLight}{HTML}{FFF3D7}
\definecolor{Ink}{HTML}{172A32}
\definecolor{Grey}{HTML}{64757D}
\definecolor{White}{HTML}{FFFFFF}

\pagestyle{empty}
\setlength{\parindent}{0pt}
\setlength{\parskip}{0pt}
\setlength{\tabcolsep}{3pt}

% ------------------------------------------------------------
% CADRE
% ------------------------------------------------------------
\newcommand{\PageFrame}{%
\begin{tikzpicture}[remember picture,overlay]
  \coordinate (A) at ([xshift=8mm,yshift=-8mm]current page.north west);
  \coordinate (B) at ([xshift=-8mm,yshift=-8mm]current page.north east);
  \coordinate (C) at ([xshift=-8mm,yshift=8mm]current page.south east);
  \coordinate (D) at ([xshift=8mm,yshift=8mm]current page.south west);
  \def\r{1.7mm}

  \draw[Border,line width=.65pt]
    (A)
    -- ([xshift=-\r]B)
    arc[start angle=90,end angle=0,radius=\r]
    -- ([yshift=\r]C)
    arc[start angle=0,end angle=-90,radius=\r]
    -- (D) -- cycle;

  \draw[Orange,line width=3.2mm] (A) -- (D);
\end{tikzpicture}%
}

% ------------------------------------------------------------
% BOÎTES
% ------------------------------------------------------------
\tcbset{
  enhanced,
  boxrule=.6pt,
  arc=2.8mm,
  left=3.7mm,right=3.7mm,top=1.75mm,bottom=1.75mm,
  boxsep=1mm,
  coltext=Ink
}

\newtcolorbox{bluebox}{colback=OrangeLight,colframe=OrangeLine}
\newtcolorbox{redbox}{colback=RedLight,colframe=Border}
\newtcolorbox{whitebox}{colback=White,colframe=Border}
\newtcolorbox{goldbox}{
  colback=GoldLight,colframe=GoldLight,
  boxrule=0pt,arc=2.5mm,
  left=2.7mm,right=2.7mm,top=1.25mm,bottom=1.25mm
}

\newcommand{\sectiontitle}[2]{%
  {\large\bfseries\textcolor{#1}{#2}}%
}

\newcommand{\pill}[2][OrangeLight]{%
  \tikz[baseline=(X.base)]{
    \node[
      fill=#1,rounded corners=2pt,
      inner xsep=4.2pt,inner ysep=1.6pt,
      font=\small\bfseries,text=Ink
    ] (X) {#2};
  }%
}

\newcommand{\signalpill}[1]{%
  \tikz[baseline=(X.base)]{
    \node[
      fill=GoldLight,rounded corners=2pt,
      inner xsep=3.7pt,inner ysep=1.1pt,
      font=\scriptsize\bfseries,text=Ink
    ] (X) {#1};
  }%
}

% ------------------------------------------------------------
% TIMELINE
% ------------------------------------------------------------
\newcommand{\Timeline}{%
\begin{tikzpicture}[x=1cm,y=1cm]
  \draw[-{Latex[length=3mm]},very thick,Ink] (0,0) -- (15.3,0);

  \draw[-{Latex[length=3mm,width=2mm]},very thick,Orange]
    (8.2,1.10) -- (8.2,.10);
  \node[font=\bfseries\small,text=Orange] at (8.2,1.32) {MAINTENANT};

  \draw[Orange,line width=2.2pt] (9.5,0) -- (13.2,0);
  \fill[Orange] (9.5,0) circle (3.2pt);
  \fill[Orange] (11.3,0) circle (3.2pt);
  \fill[Orange] (13.2,0) circle (3.2pt);

  \node[
    fill=OrangeLight,rounded corners=2pt,
    inner xsep=5pt,inner ysep=2pt,
    font=\scriptsize\bfseries,text=OrangeDark
  ] at (11.3,.72) {actions à venir};

  \node[font=\scriptsize\bfseries,text=Grey] at (11.3,-.55)
    {projetées dans l'avenir};

  \node[font=\scriptsize\bfseries,text=Grey] at (3.0,-.48) {PASSÉ};
  \node[font=\scriptsize\bfseries,text=Orange] at (8.2,-.48) {PRÉSENT};
  \node[font=\scriptsize\bfseries,text=Grey] at (13.7,-.48) {FUTUR};
\end{tikzpicture}%
}

\begin{document}
\PageFrame

% ------------------------------------------------------------
% TITRE
% ------------------------------------------------------------
\vspace*{1mm}
\begin{center}
  {\fontsize{25}{27}\selectfont\bfseries\textcolor{Ink}{FUTUR SIMPLE}}

  \vspace{0.75mm}
  {\large\textcolor{Grey}{Le futur simple}}

  \vspace{3mm}

  \begin{tcolorbox}[
    width=.88\textwidth,
    colback=OrangeLight,
    colframe=OrangeLight,
    boxrule=0pt,
    arc=3mm,
    halign=center,
    top=2.2mm,bottom=2.2mm
  ]
    {\large\bfseries\textcolor{OrangeDark}{INFINITIF + TERMINAISONS DU FUTUR}}
  \end{tcolorbox}
\end{center}

\vspace{1mm}

% ------------------------------------------------------------
% TIMELINE
% ------------------------------------------------------------
\begin{bluebox}
  \sectiontitle{Orange}{La ligne du temps}

  \vspace{0.75mm}
  \begin{center}
    \Timeline
  \end{center}

  \vspace{-1mm}
  \begin{center}
    \small
    Le \textbf{futur simple} présente des actions
    \textbf{\textcolor{Orange}{situées dans l'avenir, après le moment présent}}.
  \end{center}
\end{bluebox}

\vspace{1.35mm}

% ------------------------------------------------------------
% CONSTRUCTION
% ------------------------------------------------------------
\begin{bluebox}
  \sectiontitle{Orange}{Construction}

  \vspace{0.75mm}
  \begin{tabularx}{\textwidth}{@{}p{.29\textwidth} X@{}}
    \textbf{Avec un verbe régulier} &
      Je \textbf{parlerai}. \qquad Tu \textbf{finiras}. \qquad Ils \textbf{vendront}.\\
    \textbf{Terminaisons} &
      \textbf{-ai, -as, -a, -ons, -ez, -ont}.\\
    \textbf{Radicaux irréguliers} &
      J'\textbf{aurai}, je \textbf{serai}, je \textbf{ferai}, je \textbf{viendrai}.
  \end{tabularx}

  \vspace{0.7mm}
  \textit{Pour les verbes réguliers, les terminaisons s'ajoutent à l'infinitif. Certains verbes ont un radical irrégulier.}
\end{bluebox}

\vspace{1.35mm}

% ------------------------------------------------------------
% EXEMPLES / ATTENTION
% ------------------------------------------------------------
\begin{minipage}[t]{.485\textwidth}
\begin{bluebox}
  \sectiontitle{Orange}{Exemples}

  \vspace{0.75mm}
  Demain, je \textbf{partirai} tôt.\\[-1mm]
  Nous \textbf{visiterons} le village samedi.\\[-1mm]
  Elle \textbf{finira} son travail ce soir.
\end{bluebox}
\end{minipage}
\hfill
\begin{minipage}[t]{.485\textwidth}
\begin{redbox}
  \sectiontitle{Red}{Attention !}

  \vspace{0.75mm}
  Ne pas confondre le futur simple
  avec le \textbf{conditionnel} :
  \textbf{je parlerai} / \textbf{je parlerais}.

  \vspace{0.4mm}
  À l'écrit, l'accent ou la terminaison
  change le sens et la prononciation.
\end{redbox}
\end{minipage}

\vspace{1.35mm}

% ------------------------------------------------------------
% USES
% ------------------------------------------------------------
\begin{bluebox}
  \sectiontitle{Orange}{Quand utilise-t-on le futur simple ?}

  \vspace{0.75mm}
  \begin{tabularx}{\textwidth}{@{}p{.31\textwidth} X@{}}
    \pill{AVENIR}
      &
      Action ou événement situé après le moment présent.
      \newline \textit{Je partirai demain matin.}\\[0.7mm]

    \pill{PRÉVISION}
      &
      Ce que l'on pense ou prévoit pour l'avenir.
      \newline \textit{Il fera beau demain.}\\[0.7mm]

    \pill{PROMESSE / ENGAGEMENT}
      &
      Promesse, engagement ou décision présentée comme future.
      \newline \textit{Je t'aiderai demain.}\\[0.7mm]

    \pill{FUTUR CERTAIN / PROGRAMMÉ}
      &
      Fait considéré comme devant se produire.
      \newline \textit{La réunion commencera à 9 heures.}
  \end{tabularx}
\end{bluebox}

\vspace{1.35mm}

% ------------------------------------------------------------
% COMPARISON
% ------------------------------------------------------------
\begin{whitebox}
  \sectiontitle{Orange}{Futur simple ou futur proche ?}

  \vspace{0.65mm}
  \begin{tabularx}{\textwidth}{@{}p{.45\textwidth} X@{}}
    \textbf{\textcolor{Orange}{Futur simple}} &
      \textbf{\textcolor{Grey}{Futur proche}}\\
    avenir envisagé ou présenté comme futur &
      avenir proche, intention ou événement imminent\\[-.2mm]
    \textit{Je partirai demain.} &
      \textit{Je vais partir.}
  \end{tabularx}
\end{whitebox}

\vspace{1.35mm}

% ------------------------------------------------------------
% SIGNAL WORDS
% ------------------------------------------------------------
\begin{goldbox}
  \begin{center}
    {\normalsize\bfseries\textcolor{Gold}{REPÈRES FRÉQUENTS}}
    \hspace{2mm}
    \signalpill{demain}\hspace{1mm}
    \signalpill{bientôt}\hspace{1mm}
    \signalpill{plus tard}

    \vspace{0.75mm}

    \signalpill{la semaine prochaine}\hspace{1mm}
    \signalpill{dans deux jours}\hspace{1mm}
    \signalpill{l'année prochaine}
  \end{center}
\end{goldbox}

\end{document}
